% GENERATED FILE. Edit canonical lessons, then run npm run generate:books.
% canonical-source-hash: fnv1a64:1b8fd54bad6e78b2
% canonical-lessons: HI-C481-letna, HI-C481-sikorna, HI-C481-lagrana, HI-C481-ucakna, HI-C481-kudna

\chapter{Lying Down and Getting About}
\label{ch:hi-lying-down-and-getting-about}
\hlchaptermodality{481}

\begin{chapteropening}
\textbf{By the end of this chapter:} \emph{I can say how someone moved, slipped or fell.}

Complete the last lesson of chapter 481: I can say how someone moved, slipped or fell.
\end{chapteropening}

\section[\texorpdfstring{\hi{लेटना} (leṭnā)}{लेटना (leṭnā)}]{\hi{लेटना} (leṭnā) — to lie down}

\label{lesson:HI-C481-letna}



\begin{quote}
Before the new one: say the Hindi for an aroma, then the Hindi for to lick.
\end{quote}

\subsection*{You'll want to know: \hi{लेटना}}
\textbf{\hi{लेटना}} — \emph{leṭnā} — \textquotedblleft{}to lie down\textquotedblright{}.

\begin{quote}
\hi{खाने} \hi{के} \hi{बाद} \hi{दादा} \hi{थोड़ी} \hi{देर} \hi{लेटते} \hi{हैं।}
\end{quote}

\begin{quote}
\emph{khāne ke bād dādā thoṛī der leṭte haiṁ.}
\end{quote}

\begin{quote}
After eating, my grandfather lies down for a while.
\end{quote}

Past (he …): \hi{वह} \hi{बिस्तर} \hi{पर} \hi{लेट} \hi{गया}, \emph{vah bistar par leṭ gayā}.

\begin{cousinweb}
\textbf{\hi{लेटना}} (\emph{leṭnā}) is a tadbhava: CDIAL lists Hindi \emph{leṭnā}, \textquotedblleft{}to lie down\textquotedblright{}, under Sanskrit \emph{lēṭyáti} (CDIAL 11109).
\end{cousinweb}

\begin{grammarlens}[title={What you've built}]
One more word for describing what you lived and felt.
\end{grammarlens}

\subsection*{Your turn}
\begin{itemize}
\raggedright
  \item \emph{Say it:} \emph{leṭnā}
  \item \emph{Say it:} \emph{leṭnā}, once more
  \item \emph{Say it:} \emph{khāne ke bād dādā thoṛī der leṭte haiṁ.} — After eating, my grandfather lies down for a while.
  \item \emph{Recall:} say the Hindi for an aroma, then the Hindi for to lick, then say \emph{leṭnā} again
\end{itemize}

\begin{tcolorbox}[breakable,colback=teal!4,colframe=teal!35!black,title={Before you move on}]
What does \hi{लेटना} mean? (\textquotedblleft{}To lie down\textquotedblright{}.) Say it once more.
\end{tcolorbox}
\section[\texorpdfstring{\hi{सिकोड़ना} (sikoṛnā)}{सिकोड़ना (sikoṛnā)}]{\hi{सिकोड़ना} (sikoṛnā) — to draw in; to screw up (the face)}

\label{lesson:HI-C481-sikorna}



\begin{quote}
Before the new one: say the Hindi for to lick, then the Hindi for to lie down.
\end{quote}

\subsection*{You'll want to know: \hi{सिकोड़ना}}
\textbf{\hi{सिकोड़ना}} — \emph{sikoṛnā} — \textquotedblleft{}to draw in; to screw up (the face)\textquotedblright{}.

\begin{quote}
\hi{ठंड} \hi{में} \hi{दादी} \hi{ने} \hi{रजाई} \hi{में} \hi{पैर} \hi{सिकोड़े।}
\end{quote}

\begin{quote}
\emph{ṭhaṇḍ meṁ dādī ne rajāī meṁ pair sikoṛe.}
\end{quote}

\begin{quote}
In the cold, my grandmother drew her feet up under the quilt.
\end{quote}

Past, with \hi{ने} (he or she …): \hi{उसने} \hi{नाक} \hi{सिकोड़ी}, \emph{usne nāk sikoṛī}.

\begin{cousinweb}
\textbf{\hi{सिकोड़ना}} (\emph{sikoṛnā}) is a tadbhava: CDIAL lists Hindi \emph{sikoṛnā}, \textquotedblleft{}to contract, compress, crumple together\textquotedblright{}, under \emph{*saṁkōṭayati}, a form it reconstructs (CDIAL 12833). \emph{nāk sikoṛnā}, \textquotedblleft{}to screw up one's nose\textquotedblright{}, is to show distaste.
\end{cousinweb}

\begin{grammarlens}[title={What you've built}]
One more word for describing what you lived and felt.
\end{grammarlens}

\subsection*{Your turn}
\begin{itemize}
\raggedright
  \item \emph{Say it:} \emph{sikoṛnā}
  \item \emph{Say it:} \emph{sikoṛnā}, once more
  \item \emph{Say it:} \emph{ṭhaṇḍ meṁ dādī ne rajāī meṁ pair sikoṛe.} — In the cold, my grandmother drew her feet up under the quilt.
  \item \emph{Recall:} say the Hindi for to lick, then the Hindi for to lie down, then say \emph{sikoṛnā} again
\end{itemize}

\begin{tcolorbox}[breakable,colback=teal!4,colframe=teal!35!black,title={Before you move on}]
What does \hi{सिकोड़ना} mean? (\textquotedblleft{}To draw in; to screw up (the face)\textquotedblright{}.) Say it once more.
\end{tcolorbox}
\section[\texorpdfstring{\hi{लँगड़ाना} (lãgṛānā)}{लँगड़ाना (lãgṛānā)}]{\hi{लँगड़ाना} (lãgṛānā) — to limp}

\label{lesson:HI-C481-lagrana}



\begin{quote}
Before the new one: say the Hindi for to lie down, then the Hindi for to draw in.
\end{quote}

\subsection*{You'll want to know: \hi{लँगड़ाना}}
\textbf{\hi{लँगड़ाना}} — \emph{lãgṛānā} — \textquotedblleft{}to limp\textquotedblright{}.

\begin{quote}
\hi{मोच} \hi{की} \hi{वजह} \hi{से} \hi{अरुण} \hi{लँगड़ा} \hi{रहा} \hi{था।}
\end{quote}

\begin{quote}
\emph{moc kī vajah se Aruṇ lãgṛā rahā thā.}
\end{quote}

\begin{quote}
Arun was limping because of a sprain.
\end{quote}

Past (he …): \hi{वह} \hi{लँगड़ाया}, \emph{vah lãgṛāyā}.

\begin{cousinweb}
\textbf{\hi{लँगड़ाना}} (\emph{lãgṛānā}) is a tadbhava: CDIAL lists Hindi \emph{lãgṛānā}, \textquotedblleft{}to limp\textquotedblright{}, under \emph{*lakka}, a form it reconstructs (CDIAL 10877).
\end{cousinweb}

\begin{grammarlens}[title={What you've built}]
One more word for describing what you lived and felt.
\end{grammarlens}

\subsection*{Your turn}
\begin{itemize}
\raggedright
  \item \emph{Say it:} \emph{lãgṛānā}
  \item \emph{Say it:} \emph{lãgṛānā}, once more
  \item \emph{Say it:} \emph{moc kī vajah se Aruṇ lãgṛā rahā thā.} — Arun was limping because of a sprain.
  \item \emph{Recall:} say the Hindi for to lie down, then the Hindi for to draw in, then say \emph{lãgṛānā} again
\end{itemize}

\begin{tcolorbox}[breakable,colback=teal!4,colframe=teal!35!black,title={Before you move on}]
What does \hi{लँगड़ाना} mean? (\textquotedblleft{}To limp\textquotedblright{}.) Say it once more.
\end{tcolorbox}
\section[\texorpdfstring{\hi{उचकना} (ucaknā)}{उचकना (ucaknā)}]{\hi{उचकना} (ucaknā) — to stand on tiptoe; to spring up}

\label{lesson:HI-C481-ucakna}



\begin{quote}
Before the new one: say the Hindi for to draw in, then the Hindi for to limp.
\end{quote}

\subsection*{You'll want to know: \hi{उचकना}}
\textbf{\hi{उचकना}} — \emph{ucaknā} — \textquotedblleft{}to stand on tiptoe; to spring up\textquotedblright{}.

\begin{quote}
\hi{ऊपर} \hi{की} \hi{अलमारी} \hi{तक} \hi{पहुँचने} \hi{के} \hi{लिए} \hi{मीरा} \hi{उचकी।}
\end{quote}

\begin{quote}
\emph{ūpar kī almārī tak pahũcne ke liye Mīrā ucakī.}
\end{quote}

\begin{quote}
Mira stood on tiptoe to reach the top cupboard.
\end{quote}

Past (he …): \hi{वह} \hi{उचका}, \emph{vah ucakā}.

\begin{cousinweb}
\textbf{\hi{उचकना}} (\emph{ucaknā}) is a tadbhava: CDIAL lists Hindi \emph{ucaknā}, \textquotedblleft{}to leap up\textquotedblright{}, under Sanskrit \emph{ucca}, \textquotedblleft{}high\textquotedblright{} (CDIAL 1634): to push oneself higher.
\end{cousinweb}

\begin{grammarlens}[title={What you've built}]
One more word for describing what you lived and felt.
\end{grammarlens}

\subsection*{Your turn}
\begin{itemize}
\raggedright
  \item \emph{Say it:} \emph{ucaknā}
  \item \emph{Say it:} \emph{ucaknā}, once more
  \item \emph{Say it:} \emph{ūpar kī almārī tak pahũcne ke liye Mīrā ucakī.} — Mira stood on tiptoe to reach the top cupboard.
  \item \emph{Recall:} say the Hindi for to draw in, then the Hindi for to limp, then say \emph{ucaknā} again
\end{itemize}

\begin{tcolorbox}[breakable,colback=teal!4,colframe=teal!35!black,title={Before you move on}]
What does \hi{उचकना} mean? (\textquotedblleft{}To stand on tiptoe; to spring up\textquotedblright{}.) Say it once more.
\end{tcolorbox}
\section[\texorpdfstring{\hi{कूदना} (kūdnā)}{कूदना (kūdnā)}]{\hi{कूदना} (kūdnā) — to jump}

\label{lesson:HI-C481-kudna}



\begin{quote}
Before the new one: say the Hindi for to limp, then the Hindi for to stand on tiptoe.
\end{quote}

\subsection*{You'll want to know: \hi{कूदना}}
\textbf{\hi{कूदना}} — \emph{kūdnā} — \textquotedblleft{}to jump\textquotedblright{}.

\begin{quote}
\hi{बच्चे} \hi{ख़ुशी} \hi{से} \hi{पानी} \hi{में} \hi{कूद} \hi{पड़े।}
\end{quote}

\begin{quote}
\emph{bacce khuśī se pānī meṁ kūd paṛe.}
\end{quote}

\begin{quote}
The children jumped into the water with joy.
\end{quote}

Past (he …): \hi{वह} \hi{पानी} \hi{में} \hi{कूदा}, \emph{vah pānī meṁ kūdā}.

\begin{cousinweb}
\textbf{\hi{कूदना}} (\emph{kūdnā}) is a tadbhava: CDIAL lists Hindi \emph{kūdnā}, \textquotedblleft{}to jump\textquotedblright{}, under Sanskrit \emph{kūrdati}, \textquotedblleft{}leaps\textquotedblright{} (CDIAL 3412).
\end{cousinweb}

\begin{grammarlens}[title={What you've built}]
One more word for describing what you lived and felt.
\end{grammarlens}

\subsection*{Your turn}
\begin{itemize}
\raggedright
  \item \emph{Say it:} \emph{kūdnā}
  \item \emph{Say it:} \emph{kūdnā}, once more
  \item \emph{Say it:} \emph{bacce khuśī se pānī meṁ kūd paṛe.} — The children jumped into the water with joy.
  \item \emph{Recall:} say the Hindi for to limp, then the Hindi for to stand on tiptoe, then say \emph{kūdnā} again
\end{itemize}

\begin{tcolorbox}[breakable,colback=teal!4,colframe=teal!35!black,title={Before you move on}]
What does \hi{कूदना} mean? (\textquotedblleft{}To jump\textquotedblright{}.) Say it once more.
\end{tcolorbox}
